% !PRINTED 593
\weberpage{593}
%<<CONT
\weq{1}{\mu = \overset{\varrho}{\Pi}\,\boldsymbol{s}\!\left(\frac{2\varrho}{\mu}\right)}
there occurs a number $2\varrho$ which is congruent to $\mu\eta$ with respect to
 the modulus $\mu$. Consequently $\boldsymbol{s}(\eta)$ is a divisor of $\mu$ and
hence relatively prime to $\mathfrak{m}$, as was to be proved.

Let us now assume that in the formula (11), \S\,158
\weq{2}{\tau_0 = \overset{\xi}{\Pi}\,\boldsymbol{s}\!\left(\frac{2\alpha\xi}{\mu}\right)}
$\mathfrak{m}$ is a prime ideal of the field $\Omega$, and let us accordingly set
$\mathfrak{m} = \mathfrak{p}$, so that $\xi$ runs through a complete system of residues with
respect to the modulus $\mathfrak{p}$, with the exclusion of zero.

If we then take $\mu$ so that it is divisible only by the first power
of $\mathfrak{p}$ and is not divisible by some arbitrarily chosen other prime ideal
$\mathfrak{p}'$, then by Theorem 6.\ $\tau_0$ cannot be divisible by $\mathfrak{p}'$.
Now $\tau_0$ depends only on the ideal $\mathfrak{p}$ and is independent
of how the number $\mu$ is otherwise taken.
Consequently by Theorem 6.\ $\tau_0$ is divisible by no prime ideal
different from $\mathfrak{p}$. But neither can $\mathfrak{p}$ divide $\tau_0$ in a higher
than the first power, for the factors of the
product (1) all occur among the factors of the product (2),
and consequently $\mu$ is divisible by $\tau_0$. Accordingly we may
set outright
\weq{3}{\tau_0 = \mathfrak{p}}
and since $\tau_0$ is a number in the field $\mathfrak{L}$, it follows:

7. \gesperrt{Every odd prime ideal of the field $\Omega$ is a
field $\mathfrak{L}$, a principal ideal.}
% [sic] German source prints ``ist ein Körper $\mathfrak L$, ein Hauptideal''; sense requires ``ist im Körper $\mathfrak L$ ein Hauptideal''.

The assumption that $\mu$ should be relatively prime to $\mathfrak{a}$ likewise
involves no restriction of this theorem. For, with $\mathfrak{p}$ given, we can
select from among the equivalent forms $(a, b, c)$, without changing $\lambda$,
one for which this requirement is fulfilled.

This does not, however, yet prove that $\mathfrak{p}$ is also a principal
ideal in the class field $\mathfrak{K}(\Delta)$ itself. At least, this follows
only in the case $\Delta \equiv 1\ (\mathrm{mod}\ 8)$, where $\mathfrak{K}$ is identical with
$\mathfrak{L}$\footnote{Cf.\ on this \S\,122, where the case of the prime 2 is also settled.}.

By Theorem 3., \S\,157 the factors of the product (2)
are all associated with one another. The number of these factors is
$N(\mathfrak{p}) - 1$, and accordingly we may also set:
\weq{4}{\mathfrak{p} = \Bigl[\boldsymbol{s}\!\left(\frac{2\alpha}{\mu}\right)\Bigr]^{N(\mathfrak{p})-1}.}
